\documentclass[11pt,a4paper]{article}

\usepackage[T1]{fontenc}
\usepackage[utf8]{inputenc}
\usepackage{lmodern}
\usepackage[margin=2.7cm]{geometry}
\usepackage{amsmath,amssymb,amsthm}
\usepackage{booktabs}
\usepackage{xcolor}
\usepackage{enumitem}
\usepackage[hidelinks]{hyperref}

\newtheorem{theorem}{Theorem}[section]
\newtheorem{proposition}[theorem]{Proposition}
\newtheorem{lemma}[theorem]{Lemma}
\newtheorem{observation}[theorem]{Observation}
\theoremstyle{definition}
\newtheorem{definition}[theorem]{Definition}
\theoremstyle{remark}
\newtheorem{remark}[theorem]{Remark}

% Marks unfinished parts of the draft.
\newcommand{\todo}[1]{\textcolor{red}{[TODO: #1]}}
% Leading digit, as in the paper.
\newcommand{\lead}{\delta}
\newcommand{\oeis}[1]{\href{https://oeis.org/#1}{#1}}

\title{The infinite path of the base-10 comma child graph\\ is unique and eventually periodic}
\author{\todo{authors}}
\date{Draft, \today}

\begin{document}
\maketitle

\begin{abstract}
  Angelini, Branicky, Resta, Sloane and Wilson~\cite{comma} introduced the child graph of comma
  sequences and showed that it contains an infinite path; the lexicographically earliest one in
  base 10 is \oeis{A367620}, whose sequence of branch choices \oeis{A399179} is known for only
  30 terms. Building on the finite-graph reduction of Dougherty-Bliss and Ter-Saakov~\cite{dbts},
  who proved that all comma sequences in bases 3 to 633 are finite, we extend their decade
  automaton to the child graph. In base 10 we find that the infinite path is unique, and that its
  branch choices are periodic from the eleventh decade on, with $314$ choices per period of $924$
  decades; in particular \oeis{A399179} is given by a finite automaton, as suggested in its OEIS
  entry. Along the way we record a correction to the skip-ahead formula of~\cite{comma}, and exact
  statistics of where successor sequences end, including a closed form for their mean lifetime.
\end{abstract}

\section{Introduction}

The comma sequence \oeis{A121805} starts at $1$, and each next term $k'$ of a term $k$ is the
smallest number such that $k' - k$ equals the two-digit number formed by the last digit of $k$ and
the first digit of $k'$. It famously dies after $2137453$ terms at $99999945$. Angelini et
al.~\cite{comma} study this sequence, its variants with other starting values and bases, and the
\emph{child graph}, in which a number may continue with any valid next term rather than the
smallest. They conjecture that all comma sequences in bases $b \ge 3$ are finite and prove it for
$b = 3$.

Dougherty-Bliss and Ter-Saakov~\cite{dbts} proved this conjecture for all bases $3 \le b \le 633$.
Their method reduces the behaviour of comma sequences between consecutive ``danger intervals''
$[d\, b^k - b^2, d\, b^k)$ to a map that depends on $k$ only modulo a period $L(b)$, which yields a
finite graph $G'_b$ that has a cycle if and only if some comma sequence is infinite; they check the
absence of cycles by computer~\cite{dbtscode}. They also prove Conjecture~7.1 of~\cite{comma}.

The child graph has at least one infinite path~\cite[Theorem~5.6]{comma}; the lexicographically
earliest one in base 10 is \oeis{A367620}, starting at the root $20$. Angelini et al.\ determined
its first $30$ branch choices, recorded as \oeis{A399179}, and could not decide which of several
candidate continuations are infinite. The OEIS entry of \oeis{A399179} asks for more terms and
suggests that the sequence might be described by a finite automaton.

\paragraph{Contributions.}
We re-implemented the algorithms of~\cite{comma}, and the decade-level version of the finite graph
of~\cite{dbts}, extended to the child graph. Our main results concern the child graph in base 10.
\begin{enumerate}[label=(\roman*)]
  \item The child graph has exactly one infinite path starting at a root; it is \oeis{A367620}
        (Theorem~\ref{thm:unique}).
  \item The branch choices of this path, \oeis{A399179}, are periodic from decade $11$ on with
        period $924$ decades, $314$ choices per period (Theorem~\ref{thm:periodic}). This gives all
        terms of \oeis{A399179}; its first $30$ terms agree with the known ones.
\end{enumerate}
We also record some smaller observations:
\begin{enumerate}[label=(\roman*),resume]
  \item The skip-ahead formulas (6.3) and (6.4) of~\cite{comma} lack a factor $b$
        (Section~\ref{sec:jumps}).
  \item Where successor sequences starting at $c \cdot 10^j$ end: the dependence on $j \bmod 924$,
        the joint distribution of start digit and end point, and the exact mean lifetime
        $(b-1)(b+2) / (2(b-2))$ decades (Section~\ref{sec:ends}).
\end{enumerate}
We also reproduce, independently, the exploration of the 50 child trees in Section~11
of~\cite{comma} (Section~\ref{sec:trees}), and the finiteness of all base-10 comma sequences
from~\cite{dbts}.

Our results rely on the periodicity statement of~\cite[Proposition~8]{dbts}, applied also to the
positions of branch points and to step counts (Section~\ref{sec:table}). Apart from that, they are
computational, and the code is described in Appendix~\ref{app:software}.

\section{Preliminaries}
\label{sec:preliminaries}

We follow the definitions and notation of~\cite{comma}. Fix a base $b \ge 3$ and write $\lead(n)$
for the leading digit of $n$. A \emph{comma-child} of $k$ is a number $k'$ with
$k' - k = (k \bmod b) \cdot b + \lead(k')$, and the \emph{comma-successor} of $k$ is its smallest
child, if any. The \emph{successor graph} $G_s$ and the \emph{child graph} $G_c$ have an edge from
$k$ to its successor and to each of its children, respectively. We use the following facts, all
from~\cite{comma}.

\begin{itemize}
  \item Every number has at most one parent and at most two children~\cite[\S2, Lemma~5.3]{comma},
        so $G_c$ is a forest; its roots are the $50$ numbers (4.6) of~\cite{comma} in base 10
        (Theorem~5.5 there).
  \item \emph{Landmines}, the numbers without children, are exactly $(b-1)\cdots(b-1)\,x\,y$ with
        $1 \le x \le b-2$ and $x + y = b - 1$~\cite[Theorem~5.1]{comma}. In base 10 these are
        $10^m - r$ with $r \in \{19, 28, 37, 46, 55, 64, 73, 82\}$.
  \item \emph{Branch points} with three or more digits are $d\,(b-1)^i\,d\,(b-1-d)$,
        $1 \le d \le b-2$; their children are $d\,(b-1)^{i+2}$ and
        $(d+1) \cdot b^{i+2}$~\cite[Theorem~5.2]{comma}.
  \item \emph{Starts}: the numbers $n \ge b^2 - 1$ that are not successors are exactly
        $c \cdot b^j$ with $j \ge 2$ and $2 \le c \le b-1$~\cite[Theorem~5.4]{comma}.
\end{itemize}

\begin{definition}
  For $k \ge 0$, \emph{decade} $k$ is the interval $[b^k, b^{k+1})$ of numbers with $k + 1$ digits.
\end{definition}

A successor path therefore enters decade $k \ge 2$ either by crossing $b^k$ from below or by starting
at one of the $b - 2$ numbers $c \cdot b^k$, and the larger child of a branch point in decade $k$ is
such a start. A tree of the child graph consists of successor paths glued together at branch
points: the smaller child continues the current successor path and the larger one begins a new one.

\paragraph{The finite graph of Dougherty-Bliss and Ter-Saakov.}
In~\cite{dbts}, the point $(d, u, k)$ denotes the number $d\,b^k - u$ with $0 \le u < b^2$, the last
term of a successor path in the danger interval $[d\,b^k - b^2, d\,b^k)$. The possible values of $u$
form a set $U(b, d)$ that does not depend on $k$~\cite[Lemma~6]{dbts}; for $d = 1$,
\[
  |U(b, 1)| = \frac{(b-1)(b+2)}{2},
\]
which is $54$ in base 10. The map from $(d, u, k)$ to the next point depends on $k$ only modulo a
period $L(b)$, for $k$ large enough~\cite[Proposition~8]{dbts}; in base 10, $L(10) = 924$. The
resulting finite graph $G'_b$ on the points $(d, u, k \bmod L(b))$ has in- and out-degree at most
one~\cite[Lemma~10]{dbts}, and all base-$b$ comma sequences are finite if and only if $G'_b$ has no
cycle~\cite[Proposition~11]{dbts}. Every path of $G'_b$ begins at a start $(d, 0, \kappa)$, so there
are $(b-2) L(b)$ paths.

\section{Block jumps}
\label{sec:jumps}

Let $k = f\,s\,w\,x$ in base $b$, with leading digit $f$ and last digit $x$. As long as the leading
digit stays $f$, the successor of $k$ is $k + x b + f$, and the last digit advances by $f$ modulo
$b$. Writing $k(0) = k$ and $x_j = (x + j f) \bmod b$, we get
\begin{equation}
  \label{eq:63}
  k(n) = k + b \sum_{j=0}^{n-1} x_j + n f .
\end{equation}
Since $x_j$ is periodic in $j$ with period dividing $b$, a block of $b$ steps adds
\begin{equation}
  \label{eq:D}
  D_f(x) = b \Bigl( f + \sum_{j=0}^{b-1} \bigl( (x + j f) \bmod b \bigr) \Bigr),
\end{equation}
so that
\begin{equation}
  \label{eq:64}
  k(m b) = k + m \, D_f(x) .
\end{equation}

\begin{remark}
  Equations (6.3) and (6.4) of~\cite{comma} print the sum without the factor $b$, i.e.\
  $k(mb) = k + m(bf + \sum_j x_j)$. The worked example in Section~6 of~\cite{comma} agrees with
  \eqref{eq:64}: for $f = 1$ the comma-numbers $81, 91, 1, \dots, 71$ of a block sum to
  $460 = 10\,(1 + 45) = D_1(8)$, whereas the printed formula gives $55$. The published
  version~\cite{comma} has the same formulas. Since the lengths (1.2) reported in~\cite{comma} agree
  with the corrected formula (see below), this appears to be a typographical error only.
  \cite[Equation~(3)]{dbts} describes the same blocks by the list of comma-numbers of one cycle,
  whose sum is $D_f(x)$.
\end{remark}

In base 10, $D_f(x)$ depends on $x$ only through $x \bmod \gcd(f, 10)$; its values are
\begin{center}
  \begin{tabular}{@{}cccccccccc@{}}
    \toprule
    $f$ & 1 & 2 & 3 & 4 & 5 & 6 & 7 & 8 & 9 \\
    \midrule
    $D_f$ & 460 & 420, 520 & 480 & 440, 540 & 300, 400, \dots, 700 & 460, 560 & 520 & 480, 580 & 540 \\
    \bottomrule
  \end{tabular}
\end{center}

\paragraph{How far one may jump.}
Let $p = b^{d-1}$, where $d$ is the number of digits of $k$. The successor sequence is increasing,
so all skipped terms keep leading digit $f$ provided $k(mb) < (f+1)\,p$; a child with a smaller
leading digit is impossible, so $x b + f$ is the smallest comma-number at every skipped step. In the
child graph a skipped term may not have a second child either. Such a child has leading digit
$f + 1$ (or is in the next decade), hence is at least $(f+1)\,p$, and it is less than $b^2$ above
its parent. So it suffices that
\[
  k(mb) \le (f+1)\,p - b^2 .
\]
This matches the margin in condition (6.1) of~\cite{comma}. \todo{(6.1) reads
$k \le (f+1)(b^{i+2} - b^2)$; check whether $(f+1)\,b^{i+2} - b^2$ was intended.}

\paragraph{Validation.}
Alternating maximal jumps with single steps across leading-digit changes reproduces the lengths
(1.2) and final terms (1.3) of~\cite{comma} for the starting values $1, \dots, 8$; the longest
case, $2.1 \cdot 10^{17}$ terms from the start $6$, takes well under a second.

\section{The trees of the child graph}
\label{sec:trees}

Section~11 of~\cite{comma} reports that, of the 50 trees of the base-10 child graph, only the tree
rooted at $20$ is infinite, and that the longest-lived rival, rooted at $30$, persists until
$10^{365} - 82$. As a first check of our implementation we reproduced this by walking every tree
completely, using the child-graph jump bound of Section~\ref{sec:jumps} for unbranched stretches.
Since paths never merge, each node is visited at most once, and the cost grows with the number of
digits reached and the number of branch points rather than with the number of nodes.

\begin{proposition}[{cf.~\cite[\S11]{comma}}]
  Following all paths up to $1000$ digits, every tree of the base-10 child graph except the one
  rooted at $20$ is finite. The tree rooted at $30$ has $1007$ branch points and $1008$ leaves; its
  largest leaf is exactly $10^{365} - 82$, reached after about $2 \cdot 10^{363}$ steps from the
  root. No other tree reaches $10^{133}$.
\end{proposition}

\begin{remark}
  Section~11 of~\cite{comma} says that the root-$30$ path ``persisted for $10^{365} - 82$ terms''.
  The number $10^{365} - 82$ is the value of the landmine; the path itself has about
  $2 \cdot 10^{363}$ terms, in line with the rule of thumb that the $n$-th term is about $50 n$.
\end{remark}

Table~\ref{tab:trees} lists the trees with at least ten leaves. All 49 finite trees are listed in
Appendix~\ref{app:trees}. Trees without branch points are single successor paths, and for the roots
$3, 4, 7$ they reproduce the corresponding entries of (1.2) and (1.3).

\begin{table}[ht]
  \centering
  \begin{tabular}{@{}rrrcc@{}}
    \toprule
    root & leaves & branch points & height (edges)                   & largest leaf              \\
    \midrule
       6 &     32 &            31 & $\approx 10^{40.31\phantom{0}}$  & $10^{42\phantom{0}} - 64$ \\
      13 &     66 &            65 & $\approx 10^{85.31\phantom{0}}$  & $10^{87\phantom{0}} - 37$ \\
      17 &     31 &            30 & $\approx 10^{52.32\phantom{0}}$  & $10^{54\phantom{0}} - 55$ \\
      30 &   1008 &          1007 & $\approx 10^{363.30}$            & $10^{365}           - 82$ \\
      38 &     10 &             9 & $\approx 10^{15.32\phantom{0}}$  & $10^{17\phantom{0}} - 37$ \\
      64 &    285 &           284 & $\approx 10^{130.33}$            & $10^{132}           - 19$ \\
    \bottomrule
  \end{tabular}
  \caption{The finite trees of the base-10 child graph with at least ten leaves.}
  \label{tab:trees}
\end{table}

Every positive integer lies in exactly one tree. As the 49 finite trees contain no number above
$10^{365}$, the tree rooted at $20$ contains \emph{every} number above $10^{365}$. In particular,
above this bound its successor paths are exactly the successor sequences starting at the numbers
$c \cdot 10^j$, which motivates the next section.

\section{Start and end points of successor sequences}
\label{sec:ends}

For $2 \le c \le 9$ and $j \ge 2$, the successor sequence starting at $c \cdot 10^j$ ends at a
landmine $10^m - r$ with $r \in \{19, 28, \dots, 82\}$. We call $m - j$ its \emph{lifetime} in
decades. In the terminology of~\cite{dbts}, these sequences are the paths of $G'_{10}$, and several
of the observations below are direct consequences of the structure established there. We first made
them by direct computation for all $c$ and $366 \le j < 2274$ (about $15\,000$ sequences); they are
reproduced for all $j$ by the decade table of Section~\ref{sec:table}.

\begin{observation}[Periodicity]
  \label{obs:periodic}
  $(m - j, r)$ depends only on $c$ and $j \bmod 924$, for all $j \ge 5$. No smaller period works,
  neither for the pair nor for $m - j$ or $r$ alone. Some starts with $j \le 4$ do not follow the
  periodic pattern.
\end{observation}

The period is $L(10) = 924$ of~\cite[Proposition~8]{dbts}; the observation adds that it is sharp
and that it applies from $j = 5$ on.

\begin{observation}[Entry offsets and balance]
  \label{obs:balance}
  For $k \ge 3$, the first term $\ge 10^k$ of a successor path is $10^k + u$ with $u$ in the set
  \[
    E = \{10a : 0 \le a \le 9\} \cup \{10a + e : 0 \le a,\ a + 2 \le e \le 9\}
  \]
  of $46$ offsets. In every decade $k \ge 3$, the $46$ entries and the $8$ starts are mapped
  one-to-one onto the $46$ exits into decade $k + 1$ and the $8$ landmines of the decade.
\end{observation}

The $46 + 8 = 54$ outcomes correspond to the $|U(10, 1)| = 54$ values of $u$ in~\cite{dbts}: the
last term below $10^{k+1}$ is either a landmine or the predecessor of an exit. The balance is the
decade-level form of the fact that $G'_{10}$ is a disjoint union of paths starting at the starts;
the implementation of~\cite{dbts} checks finiteness by exactly this count.

Every landmine is the end of exactly one successor path, starting either at some $c \cdot 10^j$ or
at one of the $54$ non-successors below $100$ (the list (4.1) of~\cite{comma}). The latter all die
below $10^{23}$, the last being the sequence from $40$, which ends at $10^{23} - 46$. So every
landmine above $10^{23}$ is the end of exactly one successor sequence starting at some
$c \cdot 10^j$.

\paragraph{Mean lifetime.}
Counting decade boundaries over one period of $G'_b$ gives the mean lifetime exactly: each of the
$(b-2) L(b)$ paths starting in one period crosses as many decade boundaries as its lifetime, and each
of the $|U(b, 1)|\, L(b)$ boundary points lies on exactly one path. Hence, in any base in which all
comma sequences are finite,
\[
  \text{mean lifetime} = \frac{|U(b, 1)|}{b - 2} = \frac{(b-1)(b+2)}{2(b-2)} = \frac{b+3}{2} + \frac{2}{b-2}
\]
decades, averaged over the starts $c \cdot b^j$ of one period. In base 10 this is $27/4$; we
confirmed the formula numerically for $3 \le b \le 10$. The lifetimes decay roughly geometrically,
by a factor close to $1 - 8/54$ per decade in base 10, and the longest lifetime is $45$ decades.

\begin{remark}
  \cite[Conjecture~17]{dbts} states that the path length $|P|$ in $G'_b$ is approximately
  exponentially distributed with mean close to $b/2 + 1$. If $|P|$ is measured in decade
  boundaries and averaged over paths, the formula above gives its mean exactly, which exceeds
  $b/2 + 1$ by $1/2 + 2/(b-2)$. \todo{Check that this reading of $|P|$ matches~\cite{dbts}, whose
  Section~5 averages over vertices rather than paths.}
\end{remark}

\paragraph{Where sequences end.}
Table~\ref{tab:cr} shows how the end residue $r$ depends on the start digit $c$ over one period.
The distribution is close to uniform, with one conspicuous exception, $c = 7$, $r = 28$.

\begin{table}[ht]
  \centering
  \begin{tabular}{@{}c*{8}{r}@{}}
    \toprule
    $c \setminus r$ & 19 & 28 & 37 & 46 & 55 & 64 & 73 & 82 \\
    \midrule
    2 & 120 & 107 & 110 & 109 & 126 & 136 & 102 & 114 \\
    3 & 129 &  97 & 125 & 111 & 118 & 123 & 118 & 103 \\
    4 & 100 & 105 & 123 & 123 & 128 & 101 & 129 & 115 \\
    5 & 136 &  94 & 106 & 119 & 109 & 111 & 122 & 127 \\
    6 & 123 &  75 & 116 & 116 & 120 & 113 & 135 & 126 \\
    7 &  96 & 234 & 119 & 114 &  87 &  91 &  96 &  87 \\
    8 & 120 &  87 & 114 & 119 & 133 & 115 & 106 & 130 \\
    9 & 100 & 125 & 111 & 113 & 103 & 134 & 116 & 122 \\
    \bottomrule
  \end{tabular}
  \caption{Number of $j$ in one period ($924$ consecutive values) for which the sequence starting at
  $c \cdot 10^j$ ends at a landmine $10^m - r$.}
  \label{tab:cr}
\end{table}

\begin{observation}
  For every $j \ge 5$ with $j \equiv 5 \pmod 6$, the sequence starting at $7 \cdot 10^j$ ends at
  $10^{j+1} - 28 = 99\cdots972$, within its own decade.
\end{observation}

This accounts for $154 = 924 / 6$ of the $234$ cases of $7 \mapsto 28$. The period $6$ is that of
$10^j$ modulo $13$, which divides $D_7 = 520$; a sequence starting at $7 \cdot 10^j$ passes only
through stretches with $D_7 = 520$, $D_8 \in \{480, 580\}$ and $D_9 = 540$ before its first
landmine. \todo{Turn this into a proof, and explain why $6 \mapsto 28$ is correspondingly rare.}

\section{The decade table}
\label{sec:table}

Our computations use the finite graph of~\cite{dbts} at the granularity of decades rather than
danger intervals, extended by the information needed for the child graph.

\begin{definition}
  Let $k \ge 2$. The \emph{inputs} of decade $k$ are the entries $b^k + u$ with $u \in E$ and the
  starts $c \cdot b^k$ with $2 \le c \le b - 1$. The \emph{passage} of an input is the successor
  path from it up to and including its first term in decade $k + 1$ (the \emph{exit}
  $b^{k+1} + u'$), or up to the landmine $b^{k+1} - r$ where it dies. We record the exit or
  landmine, the leading digits $d$ of the branch points on the way, and the numbers of steps.
\end{definition}

A passage consists of at most $b - 1$ stretches of constant leading digit. Let $M$ be the least
common multiple of all block increments $D_f(x)$; in base 10,
$M = 2^5 \cdot 3^3 \cdot 5^3 \cdot 7 \cdot 11 \cdot 13 \cdot 23 \cdot 29 = 72\,108\,036\,000$.

\begin{lemma}
  \label{lem:reduction}
  Let $k, k' \ge 3$ with $b^k \equiv b^{k'} \pmod M$. Then every input has the same exit or landmine
  in decade $k$ as in decade $k'$, and passes branch points with the same leading digits in the
  same order. The numbers of steps of its passages differ by $b\,(b^k - b^{k'}) \sum 1/D$, the sum
  running over the stretches of the passage, with $D$ the block increment of each stretch.
\end{lemma}

\begin{proof}
  Write $P = b^k$ and a term of stretch $f$ of decade $k$ as $n = fP + o$ with $0 \le o < P$, and
  let $t = P - o$ be its distance to the end of the stretch.

  \emph{The rule near a boundary.} As $b \mid P$, the last digit of $n$ is $x = (-t) \bmod b$. A
  child $n + xb + e$ is less than $n + b^2 \le n + P$, so its leading digit is $f$ if $xb + e < t$,
  and $f + 1$ otherwise, or $1$ in decade $k + 1$ if $f = b - 1$. Hence $n$ has the child with
  $e = f$ if and only if $xb + f < t$, the child with $e = f + 1$ (or $e = 1$ if $f = b - 1$) if
  and only if $xb + e \ge t$, and no other child. Which children exist, and where they lie relative
  to the boundary, therefore depends only on $f$ and $t$. In particular, a term with $t > b^2$ has
  the single child $n + xb + f$, and branch points and landmines have $t \le b^2$.

  \emph{The middle of a stretch.} A passage enters each stretch at an offset $o_0 < b^2$: the
  inputs have $o_0 = u \in E$ or $o_0 = 0$, and a child lies less than $b^2$ above its parent.
  So $t_0 = P - o_0 > b^3 - b^2 \ge b^2$. Let $x_0$ be the last digit on entry and $D = D_f(x_0)$.
  As long as the distance exceeds $b^2$, the terms are $n_0 + \sigma_j$, where $\sigma_j$ are the
  partial sums of the cycle of comma-numbers, with $\sigma_{j + b} = \sigma_j + D$. Let
  $q = \lfloor (t_0 - b^2 - 1) / D \rfloor \ge 0$. All terms up to $n_0 + qD$ have distance at
  least $t_0 - qD > b^2$, so the path reaches $n_0 + qD$ after $qb$ steps, with last digit $x_0$ and
  distance $t_0 - qD = b^2 + 1 + \bigl((t_0 - b^2 - 1) \bmod D\bigr)$. By the rule near a boundary,
  everything that follows in the stretch --- the remaining steps, a branch point or landmine on the
  way, and the offset at which the path enters the next stretch or decade --- depends only on $f$,
  $x_0$ and this distance, hence on $f$, $x_0$ and $t_0 \bmod D$.

  \emph{Comparing decades.} In decades $k$ and $k'$ an input enters its first stretch at the same
  offset, so the distances differ by $b^k - b^{k'}$, a multiple of $D$. By induction over the
  stretches, the passages agree stretch by stretch, and in each stretch $q$ differs by
  $(b^k - b^{k'}) / D$, that is, the number of steps by $b\,(b^k - b^{k'}) / D$.
\end{proof}

The second part of the proof is the argument of~\cite[Proposition~8]{dbts}, where the distance
$b^k + u$ is reduced modulo the sum of one cycle, $D_f(x)$ or a divisor of it. The lemma adds what we
need beyond their statement: an explicit range of $k$, the branch points, and the step counts. The
branch points could also be read off their points $(d, u, k)$: since the smaller child of the branch
point below $(d+1)\,b^k$ is $(d+1)\,b^k - 1$~\cite[Theorem~5.2]{comma}, a successor path passes it if
and only if $(d+1)\,b^k - 1$ is its last term below $(d+1)\,b^k$, i.e.\ $u = 1$. The step counts are
used only to compute term indices, not in the results of Section~\ref{sec:consequences}.

\begin{corollary}
  \label{cor:period}
  In base 10, the passages through decades $k$ and $k + 924$ agree, in the sense of
  Lemma~\ref{lem:reduction}, for all $k \ge 5$.
\end{corollary}

\begin{proof}
  For $k \ge 5$, $10^k$ is divisible by $2^5 \cdot 5^3$, and modulo the remaining factor
  $3^3 \cdot 7 \cdot 11 \cdot 13 \cdot 23 \cdot 29$ of $M$ it has period $924$, the lcm of the
  orders $3, 6, 2, 6, 22, 28$ of $10$ modulo $27, 7, 11, 13, 23, 29$.
\end{proof}

The threshold is sharp: some passages through decades $3$ and $4$ differ from those through
decades $927$ and $928$. The period $924$ is $L(10)$ of~\cite{dbts}.

\paragraph{Implementation.}
The proof of Lemma~\ref{lem:reduction} uses $b^k$ only through its residue modulo $M$, the fact that
it is a multiple of $b$, and $b^k - b^2 \ge b^2$. So it applies verbatim if $b^k$ is replaced by a
number $B \equiv b^k \pmod M$ with $B \ge 2b^2$, writing the numbers of a decade as $f B + o$, and
the implementation computes with such representatives; the step counts for $b^k$ follow from those
for $B$ as in the lemma, which makes term indices exact. Its choice of representatives repeats from
$k = 11$ on, so the table consists of the passages through the decades $11, \dots, 934$, one for each
of the $54$ inputs, $924 \cdot 54 = 49\,896$ in total; this is also the number of vertices counted by
the program of~\cite{dbtscode} for $b = 10$. By Corollary~\ref{cor:period}, the table also describes
the decades $5, \dots, 10$.

\paragraph{Validation.}
\begin{itemize}
  \item For bases $3, 4, 5, 6$ and $10$, the passages computed with the reduced $B$ agree with those
        computed with the real $b^k$ (exit, landmine, branch points and step counts) for the
        first $30$ decades of the periodic range and for $k = 400, 401, 402$.
  \item Passages through small decades agree with step-by-step computation, including the positions
        of all branch points.
  \item In base 3 the table has period $4$ and entry offsets $\{0, 2, 3, 6\}$, and it reproduces the
        transition table in the proof of Theorem~9.1 of~\cite{comma}.
\end{itemize}

\section{The child graph in base 10}
\label{sec:consequences}

\subsection{Successor sequences}

Following each entry of the table until it dies, we confirm the result of~\cite{dbts} for $b = 10$:
the successor automaton has no cycle, so every base-10 comma sequence is finite. In addition, every
successor path that enters a decade $k \ge 11$ dies within $44$ decades, and every sequence starting
at $c \cdot 10^j$ with $j \ge 11$ within $45$ decades. In base 3 the same computation recovers
Theorem~9.1 of~\cite{comma}, with a maximum lifetime of $6$ decades.

\subsection{The infinite path is unique}

In the child graph, a path may switch at a branch point in decade $k$ with leading digit $d$ to the
start $(d + 1) \cdot 10^k$. The set of exits reachable from an input is therefore also given by the
table, and the entries from which an infinite path continues are the greatest fixed point of the
condition ``some reachable exit is again such an entry''.

\begin{theorem}
  \label{thm:unique}
  For every decade $k \ge 11$ there is exactly one offset $u_k \in E$ such that an infinite path of
  the base-10 child graph passes through $10^k + u_k$. Consequently the child graph has exactly one
  infinite path starting at a root, namely \oeis{A367620}, which starts at $20$.
\end{theorem}

\begin{proof}
  The first statement is a finite computation on the table. Every infinite path meets every decade
  from some point on, so it passes through $10^k + u_k$ for all large $k$. Since every node has a
  unique parent, two paths through the same node share all earlier nodes; hence all infinite paths
  from roots coincide.
\end{proof}

All $46$ offsets occur as $u_k$, each for between $12$ and $26$ of the $924$ decades of a period.
Theorem~\ref{thm:unique} answers the question at the end of Section~11 of~\cite{comma}, where
several candidate continuations remained, and is the base-10 analogue of Theorem~10.1
of~\cite{comma} for base 3; our computation reproduces that theorem, including the alternation
between the smaller and the larger child.

\subsection{The branch choices are periodic}

Following \oeis{A367620} decade by decade, at each branch point we take the smaller child unless no
infinite path continues from it. This decision depends only on the state $(k \bmod 924, u_k)$, so:

\begin{theorem}
  \label{thm:periodic}
  The branch choices \oeis{A399179} of \oeis{A367620} are periodic from decade $11$ on, with period
  $924$ decades. One period contains $314$ branch points, at $171$ of which the path takes the
  larger child. Before decade $11$ there is a single branch point, $1\,9^8\,18 = 19999999918$ in
  decade $10$, where the path takes the smaller child.
\end{theorem}

In particular, \oeis{A399179} is given by a finite automaton, as suggested in its OEIS entry, and
all its terms can be listed: the first term, followed by the $314$ terms of Appendix~\ref{app:choices}
repeated forever. The first $30$ terms agree with the known ones,
\[
  \texttt{001110011101100001100101111110},
\]
and we reach the next branch point $2\,9^{84}\,27$ at term $\approx 10^{84.79}$, matching the
``about $10^{84.8}$'' of~\cite[\S11]{comma}.

\begin{remark}
  Counting the root $20$ as term $1$, we find the first branch point $19999999918$ at term
  $412\,987\,859$, both via the table and by direct step-by-step counting. Section~11
  of~\cite{comma} and \oeis{A367620} (offset $1$) give $a(412\,987\,860) = 19999999918$.
  \todo{Resolve this off-by-one.}
\end{remark}

\section{Open questions and next steps}
\label{sec:open}

\begin{itemize}
  \item \textbf{Other bases.} The construction is base-independent. Determine, for the bases covered
        by~\cite{dbts}, whether the infinite path of the child graph is unique and whether its
        branch choices are eventually periodic. The periods $L(b)$ grow quickly ($48\,720$ for
        $b = 11$), so this needs a faster implementation, for example one based on~\cite{dbtscode}.
  \item \textbf{A proof without the computer.} Explain the unique immortal entry per decade, and the
        pattern of choices, from the digit dynamics, as \cite[Theorem~10.1]{comma} does in base 3.
  \item \textbf{Structure of the table.} Explain the set $E$ of entry offsets and the
        non-uniformities in Table~\ref{tab:cr} directly.
  \item \textbf{Lengths.} The table gives the exact number of terms of every passage as an affine
        function of $b^k$. Use it to study the lengths (1.2) of~\cite{comma} for all starts.
  \item \textbf{OEIS.} The terms of \oeis{A399179}, and the uniqueness of the infinite path, could be
        contributed to the OEIS.
\end{itemize}


\appendix
\section{Software and reproducibility}
\label{app:software}

All computations use the Python package in the repository~\cite{repo}. The numbers and tables in
this report were produced at commit {\footnotesize\href{https://github.com/uhrm/kangaroo-sequence/tree/1192116ae9ff5b3c9f32cd87c538cb3ba24d58a2}{\texttt{1192116ae9ff5b3c9f32cd87c538cb3ba24d58a2}}},
where the outputs of all scripts are stored in \texttt{analysis/output/}. The main modules are
\begin{itemize}
  \item \texttt{comma.py}: children, successors, and the block jumps of Section~\ref{sec:jumps};
  \item \texttt{child\_graph.py}: roots and complete exploration of child trees
        (Section~\ref{sec:trees});
  \item \texttt{decades.py}: the decade table, immortal entries and the earliest infinite path
        (Sections~\ref{sec:table} and~\ref{sec:consequences}).
\end{itemize}
Table~\ref{tab:reproduce} lists which script reproduces which result. Each script is run with
\texttt{uv run python} followed by the command shown. The test suite (\texttt{uv run pytest})
checks the results against~\cite{comma} and the decade table against direct computation, as
described in the validation paragraphs above.

\begin{table}[ht]
  \centering
  \footnotesize
  \begin{tabular}{@{}ll@{}}
    \toprule
    Result & Command (output in \texttt{analysis/output/}) \\
    \midrule
    Section~\ref{sec:trees}, Appendix~\ref{app:trees}
      & \texttt{-m kangaroo\_sequence.child\_graph} (\texttt{child\_trees.txt}) \\
    Section~\ref{sec:jumps} ($D_f$), Section~\ref{sec:table} ($M$, period)
      & \texttt{analysis/decade\_period.py} (\texttt{decade\_period.txt}) \\
    Section~\ref{sec:ends}, Table~\ref{tab:cr}
      & \texttt{analysis/successor\_ends.py} (\texttt{successor\_ends.txt}) \\
    Mean lifetime for $3 \le b \le 10$
      & \texttt{analysis/mean\_lifetime.py} (\texttt{mean\_lifetime.txt}) \\
    Section~\ref{sec:consequences}, Appendix~\ref{app:choices}
      & \texttt{analysis/infinite\_path.py 100} (\texttt{infinite\_path.txt}) \\
    The decade table
      & \texttt{analysis/export\_table.py} (\texttt{decade\_table\_base10.csv}) \\
    Remark~\ref{rem:index}
      & \texttt{analysis/first\_branch\_point.py} (\texttt{first\_branch\_point.txt}) \\
    Remark~\ref{rem:index}, step by step
      & \texttt{analysis/first\_branch\_point.c}, compiled with \texttt{gcc -O2} \\
    \bottomrule
  \end{tabular}
  \caption{Scripts reproducing the results of this report.}
  \label{tab:reproduce}
\end{table}

\section{The finite trees of the base-10 child graph}
\label{app:trees}

\begin{center}
  \small
  \begin{tabular}{@{}rrll@{\qquad}rrll@{}}
    \toprule
    root & leaves & height       & largest leaf     & root & leaves & height        & largest leaf    \\
    \midrule
       1 &      8 & $10^{23.29}$  & $10^{25} - 28$  &   31 &      1 & $1$           & $45$            \\
       2 &      3 & $10^{14.29}$  & $10^{16} - 82$  &   32 &      1 & $10^{7.30}$   & $10^{9} - 82$   \\
       3 &      1 & $1$           & $36$            &   37 &      1 & $195$         & $9918$          \\
       4 &      1 & $199899$      & $10^{7} - 55$   &   38 &     10 & $10^{15.32}$  & $10^{17} - 37$  \\
       5 &      4 & $10^{14.32}$  & $10^{16} - 64$  &   39 &      1 & $19$          & $918$           \\
       6 &     32 & $10^{40.31}$  & $10^{42} - 64$  &   40 &      4 & $10^{28.29}$  & $10^{30} - 82$  \\
       7 &      1 & $14$          & $936$           &   41 &      1 & $18$          & $954$           \\
       8 &      3 & $10^{11.30}$  & $10^{13} - 28$  &   42 &      1 & $1980$        & $99927$         \\
       9 &      3 & $10^{22.33}$  & $10^{24} - 19$  &   43 &      1 & $1$           & $81$            \\
      10 &      2 & $10^{11.29}$  & $10^{13} - 19$  &   49 &      1 & $192$         & $9945$          \\
      13 &     66 & $10^{85.31}$  & $10^{87} - 37$  &   50 &      3 & $10^{16.32}$  & $10^{18} - 73$  \\
      14 &      4 & $10^{9.30}$   & $10^{11} - 82$  &   51 &      8 & $10^{16.32}$  & $10^{18} - 55$  \\
      15 &      1 & $1$           & $72$            &   52 &      9 & $10^{19.29}$  & $10^{21} - 19$  \\
      16 &      4 & $10^{12.33}$  & $10^{14} - 55$  &   53 &      1 & $20$          & $972$           \\
      17 &     31 & $10^{52.32}$  & $10^{54} - 55$  &   54 &      1 & $0$           & $54$            \\
      18 &      1 & $0$           & $18$            &   62 &      2 & $205967$      & $10^{7} - 19$   \\
      19 &      1 & $21481$       & $999981$        &   63 &      1 & $0$           & $63$            \\
      20 &      - & -             & -               &   64 &    285 & $10^{130.33}$ & $10^{132} - 19$ \\
      21 &      3 & $2073$        & $99981$         &   65 &      1 & $2099$        & $99963$         \\
      25 &      1 & $192$         & $9963$          &   74 &      1 & $19$          & $945$           \\
      26 &      2 & $199508$      & $10^{7} - 28$   &   75 &      1 & $2135$        & $99954$         \\
      27 &      1 & $0$           & $27$            &   76 &      9 & $10^{22.29}$  & $10^{24} - 46$  \\
      28 &      1 & $196$         & $9927$          &   86 &      1 & $1982$        & $99918$         \\
      29 &      1 & $10^{6.31}$   & $10^{8} - 64$   &   87 &      2 & $10^{10.29}$  & $10^{12} - 82$  \\
      30 &   1008 & $10^{363.30}$ & $10^{365} - 82$ &   98 &      1 & $204$         & $9936$          \\
    \bottomrule
  \end{tabular}
\end{center}
Heights count edges from the root to the deepest leaf; height entries $10^{x}$ are rounded to two
decimals of the exponent. Every largest leaf is a landmine, given exactly.

\section{Branch choices of A367620 over one period}
\label{app:choices}

The first $314$ terms of \oeis{A399179}, the choices at the branch points of \oeis{A367620} in
decades $10, \dots, 933$ ($0$ = smaller child, $1$ = larger child); the sequence repeats them forever
(Theorem~\ref{thm:periodic}).
\begin{center}
  \small
  \texttt{0011 1001 1101 1000 0110 0101 1111 1000 0110 0100 1000 0110}\\
  \texttt{1010 0010 0001 1100 1011 0001 1001 1100 0111 1111 1001 0111}\\
  \texttt{0011 1110 0101 1111 1000 0011 0000 0000 1011 1110 1111 1010}\\
  \texttt{1110 1111 1110 1101 0011 0011 0001 0110 0100 1111 0001 1111}\\
  \texttt{0001 1011 1011 0010 0100 0110 1011 0111 0010 0101 1000 1100}\\
  \texttt{0011 1110 1110 0010 0010 1110 0001 1111 1101 0101 0000 1110}\\
  \texttt{0100 0111 1010 1110 1111 1010 11\phantom{xx xxxx xxxx xxxx xxxx xxxx}}
\end{center}


\bibliographystyle{plain}
\bibliography{references}

\end{document}
